% ---
% titre: Priorité des opérations avec des fractions
% theme: NO
% chapitre: Nombres rationnels
% sous_chapitre: Priorité des opérations
% niveau: [10e]
% type: EVAL
% tags: [c-fractions, c-operations, c-addition, c-soustraction, c-multiplication, c-division, p-question-TS]
% difficulte: 2
% corrige: true
% ---

\question[8]
Effectue les calculs suivants en respectant les priorités des opérations. Montre chaque étape et donne le résultat sous la forme d'une fraction simplifiée au maximum.

\vspace{10pt}
\begin{enumerate}[leftmargin=*,align=parleft, label=\alph*), itemsep=60pt]
\item $\dfrac{1}{2} + \dfrac{1}{3} \cdot \dfrac{3}{4} =$ \sol{}{$\dfrac{3}{4}$}
\item $\dfrac{5}{6} - \dfrac{2}{3} : \dfrac{4}{3} =$ \sol{}{$\dfrac{1}{3}$}
\item $\dfrac{2}{5} + \dfrac{3}{10} \cdot \dfrac{5}{6} =$ \sol{}{$\dfrac{13}{20}$}
\item $3 - \dfrac{1}{2} \cdot \dfrac{4}{3} : \dfrac{1}{6} =$ \sol{}{$\dfrac{5}{2}$}
\end{enumerate}

\begin{solution}[0pt]
\textbf{2pt par calcul :} 1pt pour la priorité respectée (la multiplication ou la division en premier), 1pt pour la suite du calcul et le résultat simplifié\\[8pt]
a) $\dfrac{1}{2} + \dfrac{1}{4} = \dfrac{2}{4} + \dfrac{1}{4} = \dfrac{3}{4}$\\[6pt]
b) $\dfrac{5}{6} - \dfrac{2}{3} \cdot \dfrac{3}{4} = \dfrac{5}{6} - \dfrac{1}{2} = \dfrac{5}{6} - \dfrac{3}{6} = \dfrac{2}{6} = \dfrac{1}{3}$\\[6pt]
c) $\dfrac{2}{5} + \dfrac{1}{4} = \dfrac{8}{20} + \dfrac{5}{20} = \dfrac{13}{20}$\\[6pt]
d) $3 - \dfrac{2}{3} : \dfrac{1}{6} = 3 - \dfrac{2}{3} \cdot \dfrac{6}{1}  = 3 - 4 = -1$
\end{solution}